\exercisesection{\S~2}

\begin{enumerate}
\def\labelenumi{\arabic{enumi})}
\tightlist
\item
  设 \(\mathfrak{g}\) 是一个 3 维单李代数，且 \(\varPhi\) 是其 Killing 型（参见~§1，习题 15、16、17）。
\end{enumerate}

\begin{enumerate}
\def\labelenumi{\alph{enumi})}
\item
  元素 \(x \in g\) 是正则元当且仅当 \(\Phi(x, x) \neq 0\)。令 \(h_x = kx\) 为由此类元素生成的嘉当子代数。证明 \(h_x\) 是分裂的当且仅当 \(2\Phi(x, x)\) 是 k 中的平方。
\item
  证明
\end{enumerate}

\(\mathfrak{g}\) 可分裂 \(\Longleftrightarrow \mathfrak{g}\) 同构于 \(\mathfrak{sl}(2,k)\) 。

\begin{enumerate}
\def\labelenumi{\arabic{enumi})}
\setcounter{enumi}{1}
\tightlist
\item
  设 \(k_{1}\) 是 \(k\) 的一个有限扩域，次数 \(n \geq 2\)。
\end{enumerate}

\begin{enumerate}
\def\labelenumi{\alph{enumi})}
\item
  证明半单 k-代数 \(\mathfrak{sl}(2,k_{1})\) 不是可分裂的。
\item
  证明可分裂的单 k-代数 \(\mathfrak{sl}(n,k)\) 包含一个非分裂的嘉当子代数 \(h_{1}\)。（选择代数 \(k_{1}\) 到 \(\mathbf{M}_{n}(k)\) 的嵌入，并取 \(\mathfrak{h}_{1}=k_{1}\cap\mathfrak{sl}(n,k)\)。）
\end{enumerate}

\begin{enumerate}
\def\labelenumi{\arabic{enumi})}
\setcounter{enumi}{2}
\tightlist
\item
  设 \((\mathfrak{g},\mathfrak{h})\) 是一个分裂半单李代数，R 是其根系，K 是 \(\mathfrak{g}\) 的 Killing 型在 \(\mathfrak{h}\) 上的限制。采用第 VI 章 §1 第 12 节的记号，有 \(K = B_{R^{\vee}}\)，且当 R 不可约时有 \(K = 4\gamma(R)\varPhi_{R^{\vee}}\)；此外，如果 R 的所有根长度相同，则
\end{enumerate}

\[
\mathrm{K} (H _ {\alpha}, H _ {\alpha}) = 4 h \quad \text { for   all } \alpha \in \mathrm{R},
\]

其中 \(h\) 是 \(\mathrm{W}(\mathbb{R})\) 的 Coxeter 数。

例如，如果 \(\mathfrak{g}\) 是 \(\mathrm{E}_6, \mathrm{E}_7\) 或 \(\mathrm{E}_8\) 型的单李代数，则 \(\mathrm{K}(H_{\alpha}, H_{\alpha})\) 等于 48、72 或 120。

\begin{enumerate}
\def\labelenumi{\arabic{enumi})}
\setcounter{enumi}{3}
\tightlist
\item
  设 \((\mathfrak{g},\mathfrak{h})\) 是一个分裂半单李代数，且 \((X_{\alpha})_{\alpha \in \mathrm{R}}\) 是满足引理 2 条件的元素族。如果 \(\alpha ,\beta \in \mathrm{R}\) 且 \(\alpha +\beta \in \mathrm{R}\)，则定义 \(\mathrm{N}_{\alpha \beta}\)，使其满足公式 \([X_{\alpha},X_{\beta}] = \mathrm{N}_{\alpha \beta}X_{\alpha +\beta}\)；如果 \(\alpha +\beta \notin \mathrm{R}\)，则置 \(\mathrm{N}_{\alpha \beta} = 0\)。证明下列公式：
\end{enumerate}

\(a)\mathrm{N}_{\alpha \beta} = -\mathrm{N}_{\beta \alpha}.\)

\begin{enumerate}
\def\labelenumi{\alph{enumi})}
\setcounter{enumi}{1}
\tightlist
\item
  如果 \(\alpha, \beta, \gamma \in \mathbb{R}\) 满足 \(\alpha + \beta + \gamma = 0\)，则
\end{enumerate}

\[
\frac {\mathrm{N} _ {\alpha \beta}}{\langle \gamma , \gamma \rangle} = \frac {\mathrm{N} _ {\beta \gamma}}{\langle \alpha , \alpha \rangle} = \frac {\mathrm{N} _ {\gamma \alpha}}{\langle \beta , \beta \rangle}.
\]

\begin{enumerate}
\def\labelenumi{\alph{enumi})}
\setcounter{enumi}{2}
\tightlist
\item
  如果 \(\alpha, \beta, \gamma, \delta \in \mathbb{R}\) 满足 \(\alpha + \beta + \gamma + \delta = 0\)，且没有一对元素的和为零，则有：
\end{enumerate}

\[
\frac {\mathrm{N} _ {\alpha \beta} \mathrm{N} _ {\gamma \delta}}{\langle \gamma + \delta , \gamma + \delta \rangle} + \frac {\mathrm{N} _ {\beta \gamma} \mathrm{N} _ {\alpha \delta}}{\langle \alpha + \delta , \alpha + \delta \rangle} + \frac {\mathrm{N} _ {\gamma \alpha} \mathrm{N} _ {\beta \delta}}{\langle \beta + \delta , \beta + \delta \rangle} = 0.
\]

\begin{enumerate}
\def\labelenumi{\arabic{enumi})}
\setcounter{enumi}{4}
\item
  设 \((\mathfrak{g},\mathfrak{h})\) 是一个分裂半单李代数，且 \((X_{\alpha})_{\alpha\in\mathbb{R}}\) 是满足引理 2 条件的元素族。令 B 为 R 的一个基。\(H_{\alpha}\)（\(\alpha\in B\)）和 \(X_{\alpha}\)（\(\alpha\in R\)）构成向量空间 g 的一个基。证明相对于此基（《代数》第 IX 章，§2），g 的 Killing 型的判别式是一个有理数，与 h、B 和 \(X_{\alpha}\) 的选择无关。¹⁰ 推出，如果 \(n=\dim g\)，则 \(\bigwedge^{n}g\) 中由 \(H_{\alpha}\)（\(\alpha\in B\)）和 \(X_{\alpha}\)（\(\alpha\in R\)）的外积定义的元素，除差一个符号外，与 h、B 和 \(X_{\alpha}\) 的选择无关。

\footnotetext[10]{关于这个判别式的显式计算，参见：T. A. SPRINGER 和 R. STEINBERG，《共轭类》（第 4.8 节），《代数群及相关有限群研讨会》，数学讲义 131，Springer-Verlag（1970）。}
  \item
  设 \((X_{\alpha})_{\alpha\in\mathbb{R}}\) 是分裂半单李代数 \((\mathfrak{g},\mathfrak{h})\) 中的一个谢瓦莱系。令 \(\alpha,\beta\in R\)，并令 p（分别为 q）是满足 \(\beta+j\alpha\in R\)（分别为 \(\beta-j\alpha\in R\)）的最大整数 j，参见~引理 4。证明
\end{enumerate}

\[
\operatorname{ad} (X _ {\alpha}) ^ {k} (X _ {\beta - q \alpha}) = \pm k! X _ {\beta + (k - q) \alpha} \text { for } 0 \leq k \leq p + q.
\]

推出命题 8 中的代数 \(\mathfrak{g}_{\mathbf{Z}}\) 在 \(\mathrm{ad}(X_{\alpha})^k / k!\) 和 \(e^{\mathrm{ad}(X_{\alpha})}\) 下稳定（参见~§12）。

\begin{enumerate}
\def\labelenumi{\arabic{enumi})}
\setcounter{enumi}{6}
\tightlist
\item
  假设 \(k\) 是有序域。设 \((\mathfrak{g},\mathfrak{h})\) 是秩为 \(l\) 的分裂半单李代数；置 \(\dim \mathfrak{g} = l + 2m\)。
\end{enumerate}

\begin{enumerate}
\def\labelenumi{\alph{enumi})}
\item
  证明 \(\varPhi\)（g 的 Killing 型）是一个秩为 2m 的中性型与一个秩为 l 的正非退化型的直和；特别地，其指标为 m。
\item
  设 \(\varphi\) 是 \(\mathfrak{g}\) 的一个对合自同构，其在 \(\mathfrak{h}\) 上的限制为 -Id。证明下列型
\end{enumerate}

\[
(x, y) \mapsto - \Phi (x, \varphi (y)), \quad x, y \in \mathfrak {g},
\]

是对称的、非退化的且正的。

\(c)\) 令 \(k' = k(\sqrt{\alpha})\)，其中 \(\alpha\) 是一个 \(< 0\) 的元素，属于 \(k\)。记 \(c\) 为 \(k\) 上 \(k'\) 的非平凡自同构。令 \(\mathfrak{g}_c\) 为 \(k\)-子空间 \(\mathfrak{g}_{(k')} = k' \otimes_k \mathfrak{g}\) 中由满足下式的元素 \(y\) 构成：

\[
(1 \otimes \varphi). y = (c \otimes 1). y.
\]

证明 \(g_{c}\) 是 \(\mathfrak{g}_{(k^{\prime})}\) 的 k-李子代数，且 \(g_{c}\) 到 \(\mathfrak{g}_{(k^{\prime})}\) 的嵌入延拓为从 \(k^{\prime}\otimes_{k}g_{c}\) 到 \(\mathfrak{g}_{(k^{\prime})}\) 的同构。代数 \(g_{c}\) 是半单的，而 \(\sqrt{\alpha}h\) 是它的一个嘉当子代数。

\(d\) ) 证明 \(\mathfrak{g}_c\) 的 Killing 型是负定的。推出 \(\mathfrak{g}_c\) 不是可分裂的（除非 \(\mathfrak{g} = 0\)）。

\begin{enumerate}
\def\labelenumi{\alph{enumi})}
\setcounter{enumi}{4}
\tightlist
\item
  当 \(k = \mathbf{R}\) 时，证明 \(\operatorname{Int}(\mathfrak{g}_c)\) 是紧的。
\end{enumerate}

\begin{enumerate}
\def\labelenumi{\arabic{enumi})}
\setcounter{enumi}{7}
\tightlist
\item
  \begin{enumerate}
  \def\labelenumii{\alph{enumii})}
  \tightlist
  \item
    设 g 是李代数，n 为一个 \(\geq 0\) 的整数。令 \(\Sigma_{n}(\mathfrak{g})\) 为 \(g^{n}\) 的子集，由作为 k-代数生成 g 的族 \((x_{1},\ldots,x_{n})\) 组成。证明 \(\Sigma_{n}(\mathfrak{g})\) 在 Zariski 拓扑中是 \(g^{n}\) 的开集（第 VII 章附录 I）。如果 \(k'\) 是 k 的扩域，则 \(\Sigma_{n}(\mathfrak{g}_{(k')})\cap\mathfrak{g}^{n}=\Sigma_{n}(\mathfrak{g})\)。推出，如果 \(\mathfrak{g}_{(k')}\) 可由 n 个元素生成，则 g 也可以。
  \end{enumerate}
\end{enumerate}

\begin{enumerate}
\def\labelenumi{\alph{enumi})}
\setcounter{enumi}{1}
\item
  设 \((\mathfrak{g},\mathfrak{h})\) 是一个分裂半单李代数。令 \(x\) 是 \(\mathfrak{h}\) 中满足：\(\alpha (x)\neq 0\) 对所有 \(\alpha \in \mathrm{R}\) 成立，且 \(\alpha (x)\neq \beta (x)\) 对任意两个不同的 \(\alpha ,\beta \in \mathrm{R}\) 成立的元素。对所有 \(\alpha \in \mathrm{R}\)，令 \(y_{\alpha}\) 为 \(\mathfrak{g}^{\alpha}\) 的一个非零元素，并令 \(y = \sum_{\alpha \in \mathrm{R}}y_{\alpha}\)。证明，对所有 \(\alpha \in \mathrm{R}\)，存在一个无常数项的多项式 \(\mathrm{P}_{\alpha}(\mathrm{T})\in k[\mathrm{T}]\)，使得 \(y_{\alpha} = \mathrm{P}_{\alpha}(\mathrm{ad}x).y\)。推出 \(\mathfrak{g}\) 由 \(\{x,y\}\) 生成。
\item
  利用 \(a\) ) 和 \(b\) )，证明每个半单李代数都可由两个元素生成。
\end{enumerate}

¶9) 设 G 是一个连通有限维实李群。设 g 为其李代数，且 \(\{x_{1},\ldots,x_{n}\}\) 是 g 的一个生成族。对 \(m\geq0\)，记 \(\Gamma_{m}\) 为由 \(\exp(2^{-m}x_{i})\)（\(1\leq i\leq n\)）生成的 G 的子群。我们有 \(\Gamma_{0}\subset\Gamma_{1}\subset\cdots\)。

\begin{enumerate}
\def\labelenumi{\alph{enumi})}
\item
  证明 \(\Gamma_{m}\) 的并在 G 中稠密。
\item
  令 \(H_{m}\) 为 \(\overline{\Gamma}_{m}\)（\(\Gamma_{m}\) 的闭包）的单位元连通分支。我们有 \(H_{0} \subset H_{1} \subset \cdots\)，且族（\(H_{m}\)）是稳定的；令 H 为 m 充分大时 \(H_{m}\) 的公共值。证明 H 在 G 中正规（注意 H 被所有 \(\Gamma_{m}\) 规范化），且 \(\Gamma_{m}\) 在 G/H 中的像是 G/H 的离散子群。这些子群的并在 G/H 中稠密，推出（参见~第 III 章，§6，习题 23 d)）G/H 是幂零的。
\item
  假设 \(\mathfrak{g} = \mathcal{D}\mathfrak{g}\)。证明 \(\mathrm{G} = \mathrm{H}\)，换言之，\(\Gamma_{m}\) 在 \(\mathrm{G}\) 中稠密，当 \(m\) 充分大时。
\end{enumerate}

\begin{enumerate}
\def\labelenumi{\arabic{enumi})}
\setcounter{enumi}{9}
\item
  设 \(G\) 是一个连通半单实李群。利用习题 8 和 9 证明，存在一个由两个元素生成的 \(G\) 的稠密子群。
\item
  设 \(\mathrm{R}\) 是秩为 \(l\) 的根系，\(\varPhi_{\mathrm{R}}\) 是相应的典则双线性型（第 VI 章，§1，第 12 节）。矩阵 \(\varPhi = (\varPhi_{\mathrm{R}}(\alpha, \beta))_{\alpha, \beta \in \mathrm{R}}\) 的秩为 \(l\)，且 \(\varPhi^2 = \varPhi\)。推出公式 \(\sum_{\alpha \in \mathrm{R}} \varPhi_{\mathrm{R}}(\alpha, \alpha) = \operatorname{Tr} \varPhi = l\)。
\end{enumerate}

